\documentclass[11pt]{article}
\usepackage[a4paper,margin=1in]{geometry}
\usepackage{fontspec}
\usepackage[boldfont=false]{xeCJK}
\setCJKmainfont{FandolSong-Regular.otf}
\usepackage{amsmath}
\usepackage{amssymb}
\usepackage{array}
\usepackage{booktabs}
\usepackage{graphicx}
\usepackage{adjustbox}
\usepackage{enumitem}
\setlist[itemize]{topsep=2pt,partopsep=0pt,parsep=0pt,itemsep=2pt,leftmargin=1.4em}
\setlist[enumerate]{topsep=2pt,partopsep=0pt,parsep=0pt,itemsep=2pt,leftmargin=1.6em}
\usepackage{parskip}
\usepackage{fvextra}
\fvset{breaklines=true,breakanywhere=true,fontsize=\small}
\setlength{\parindent}{0pt}

\begin{document}

\begin{center}
  {\LARGE\bfseries Databases}\\[0.35em]
  {\large Igcse Computer Science}
\end{center}
\vspace{0.6em}

\section*{What is a database?}

A \textbf{database} 数据库 is an organised store of data, kept so that it is easy to search, sort and update. At IGCSE you work with a \textbf{single-table database} 单表数据库 — all the data is held in one table.

\section*{Records and fields}

A database table is made of records and fields.

\begin{itemize}
\item A \textbf{record} 记录 is one row in the table — all the data about one thing (for example one student).

\item A \textbf{field} 字段 is one column in the table — one piece of data that every record has (for example "First name").

\end{itemize}

\begin{center}
\adjustbox{max width=\linewidth}{%
\begin{tabular}{lllll}
\toprule
\textbf{StudentID} & \textbf{FirstName} & \textbf{DateOfBirth} & \textbf{FormClass} & \textbf{FeesPaid} \\
\midrule
1 & Amy & 14/03/2009 & 10A & TRUE \\
2 & Ben & 02/11/2008 & 10B & FALSE \\
\bottomrule
\end{tabular}}
\end{center}

Here each \textbf{row} is a record, and each \textbf{column} is a field.

\section*{Data types for fields}

Each field stores one \textbf{data type} 数据类型. You choose the type that best fits the data.

\begin{center}
\adjustbox{max width=\linewidth}{%
\begin{tabular}{lll}
\toprule
\textbf{Data type} & \textbf{Used for} & \textbf{Example} \\
\midrule
\textbf{text/alphanumeric} 文本 & letters, digits and symbols & "10A", "Amy" \\
\textbf{character} 字符 & a single character & 'M' \\
\textbf{Boolean} 布尔值 & one of two values & TRUE / FALSE \\
\textbf{integer} 整数 & a whole number & 42 \\
\textbf{real} 实数 & a number with a decimal point & 3.5 \\
\textbf{date/time} 日期时间 & a date or a time & 14/03/2009 \\
\bottomrule
\end{tabular}}
\end{center}

\section*{The primary key}

A \textbf{primary key} 主键 is a field that holds a \textbf{unique} 唯一的 value for every record. No two records can have the same primary key, so it lets you pick out exactly one record.

In the table above, \texttt{StudentID} is a good primary key because every student has a different number. A field like \texttt{FormClass} would be a bad primary key, because many students share the same class.

\section*{Validation in a database}

When data is put into a database, \textbf{validation} 验证 checks make sure it is sensible — for example a range check on an age field, or a presence check so a field is not left empty. (You saw these checks in topic 7.)

\section*{Structured Query Language (SQL)}

\textbf{Structured Query Language} 结构化查询语言 (SQL) is a language used to \textbf{query} 查询 a database — to pick out the records you want. You must understand and complete SQL scripts.

\subsection*{SELECT, FROM and WHERE}

\begin{itemize}
\item \texttt{SELECT} says \textbf{which fields} to show.

\item \texttt{FROM} says \textbf{which table} to use.

\item \texttt{WHERE} gives a \textbf{condition} 条件, so only matching records are shown.

\end{itemize}

\begin{Verbatim}
SELECT FirstName, FormClass
FROM Student
WHERE FeesPaid = TRUE;
\end{Verbatim}

This shows the first name and class of every student who has paid the fees.

Use \texttt{*} to select \textbf{all} fields:

\begin{Verbatim}
SELECT *
FROM Student
WHERE FormClass = '10A';
\end{Verbatim}

\subsection*{ORDER BY}

\texttt{ORDER BY} sorts the results. Use \texttt{ASC} for \textbf{ascending} 升序 (smallest first, A→Z) or \texttt{DESC} for \textbf{descending} 降序 (largest first, Z→A).

\begin{Verbatim}
SELECT FirstName, DateOfBirth
FROM Student
ORDER BY DateOfBirth ASC;
\end{Verbatim}

\subsection*{AND and OR}

Join conditions with \texttt{AND} (both must be true) or \texttt{OR} (at least one must be true).

\begin{Verbatim}
SELECT FirstName
FROM Student
WHERE FormClass = '10A' AND FeesPaid = FALSE;
\end{Verbatim}

\subsection*{SUM and COUNT}

\begin{itemize}
\item \texttt{SUM} adds up the values in a number field.

\item \texttt{COUNT} counts how many records match.

\end{itemize}

\begin{Verbatim}
SELECT COUNT(StudentID)
FROM Student
WHERE FeesPaid = FALSE;
\end{Verbatim}

This counts how many students have not paid. \texttt{SUM} works the same way but adds a number field instead of counting rows.

\subsection*{Working out the output}

To find the output of an SQL script, read it in this order:

\begin{enumerate}
\item \texttt{FROM} — which table;

\item \texttt{WHERE} — keep only the records that match the condition;

\item \texttt{SELECT} — show only the chosen fields;

\item \texttt{ORDER BY} — put the results in order.

\end{enumerate}

Following these steps, you can write down exactly which rows and columns the query returns.


\end{document}
